\begin{table}[H]
\centering
\caption{Formas de $y_{h}$ en base a los tipos de raiz}
\begin{threeparttable}
{\setstretch{1.75}
\begin{tabular}{p{0.45\linewidth}p{0.45\linewidth}}\hline
    \textbf{Tipo de raiz} & \textbf{Forma de $y_{h}$} \\ \hline

    Reales y distintas ($m_{1} \neq m_{2}$) & $C_{1}x^{m_{1}x} + C_{2}x^{m_{2}x}$ \\
    Reales e iguales ($m_{1} = m_{2}$) & $x^{m} \cdot (C_{1} + C^{2} \cdot \ln{(x)})$ \\
    Complejas ($m = \alpha \pm \beta i$) & $x_{m} \cdot (\cos{(\beta \ln{(x)})} + \sin{(\beta \ln{(x)})})$ \\

    \hline
\end{tabular}
}
\begin{tablenotes} % ex: \tntote{1} ... \item[1]
    \item[] 
\end{tablenotes}
\end{threeparttable}
\label{tab:eq_cauchy_euler}
\end{table}